\chapter{Implicit Function Theorem}
\label{app:chow-iv-appendix-k}
\dcref{appk}
\source{chow-et-al-ricci-flow-part-iv:page-0348}

This appendix collects tools from functional analysis, spaces of maps,
harmonic maps, and Hodge theory that are used throughout Part IV.  All spaces
in the first section are real Banach or Hilbert spaces.

\section{The implicit function theorem}

\begin{theorem}[Lax--Milgram--Lions theorem (Theorem K.1)]
\label{thm:chow-iv-k1-lax-milgram-lions}
\source{chow-et-al-ricci-flow-part-iv:page-0348}
\dcref{appk:K.1}
\group{Implicit Function Theorem}

Let $A$ be a Hilbert space, $B$ a normed linear space, and
$f:A\times B\to\mathbb R$ a continuous bilinear form.  The following are
equivalent:
\begin{enumerate}
\item there is $c>0$ with
\[
\inf_{\|b\|=1}\sup_{\|a\|\le1}|f(a,b)|\ge c;
\]
\item for every $L\in B^*$ there is $a_L\in A$ such that
$f(a_L,b)=L(b)$ for all $b\in B$.
\end{enumerate}
When $A=B$, the condition $f(a,a)\ge c\|a\|_A^2$ implies coercivity.
\end{theorem}

\begin{lemma}[Contraction mapping theorem (Lemma K.2)]
\label{lem:chow-iv-k2-contraction}
\source{chow-et-al-ricci-flow-part-iv:page-0348}
\dcref{appk:K.2}
\group{Implicit Function Theorem}

If $(X,d)$ is complete and $f:X\to X$ satisfies
$d(f(x),f(y))\le\lambda d(x,y)$ for some $0\le\lambda<1$, then $f$ has a
unique fixed point.
\end{lemma}

\begin{theorem}[Banach inverse function theorem (Theorem K.3)]
\label{thm:chow-iv-k3-banach-inverse-function}
\source{chow-et-al-ricci-flow-part-iv:page-0349,chow-et-al-ricci-flow-part-iv:page-0350,chow-et-al-ricci-flow-part-iv:page-0351}
\uses{lem:chow-iv-k2-contraction}
\dcref{appk:K.3}
\group{Implicit Function Theorem}


Let $f:U\subset X\to Y$ be $C^k$ between Banach spaces, $k\ge1$, and let
$Df(x_0)$ be bijective.  Then on neighborhoods $W$ of $x_0$ and $f(x_0)$,
$f|_W$ is a $C^k$ diffeomorphism.  Its inverse satisfies
\[
D(f^{-1})(y)=\bigl(Df(f^{-1}(y))\bigr)^{-1}.
\]
\end{theorem}

\begin{theorem}[Banach implicit function theorem (Theorem K.4)]
\label{thm:chow-iv-k4-banach-implicit-function}
\source{chow-et-al-ricci-flow-part-iv:page-0351}
\uses{thm:chow-iv-k3-banach-inverse-function}
\dcref{appk:K.4}
\group{Implicit Function Theorem}


Let $f:\mathcal U\subset\mathcal X\times\mathcal Y\to\mathcal Z$ be $C^k$.
If $(\partial_yf)_{(x_0,y_0)}:\mathcal Y\to\mathcal Z$ is an isomorphism,
then for $x$ near $x_0$ there is a unique $C^k$ function $y=\phi(x)$ near
$y_0$ with $f(x,\phi(x))=f(x_0,y_0)$, and
\[
(D\phi)_x=-\bigl((\partial_yf)_{(x,\phi(x))}\bigr)^{-1}
\circ(\partial_xf)_{(x,\phi(x))}.
\]
\end{theorem}

\begin{definition}[Banach-manifold charts for maps (Definition K.5)]
\label{def:chow-iv-k5-map-banach-manifold}
\source{chow-et-al-ricci-flow-part-iv:page-0352,chow-et-al-ricci-flow-part-iv:page-0353}
\dcref{appk:K.5}
\group{Mapping Spaces}

For closed manifolds, the space $C^k(\mathcal M,\mathcal N)$ is charted at
$f$ by sections of $f^*T\mathcal N$ via
\[
\varphi_f(\sigma)(x)=\exp^{g_{\mathcal N}}_{f(x)}\sigma(x).
\]
For manifolds with boundary, choose a uniformly equivalent product metric with
totally geodesic boundary and restrict to sections tangent to the boundary.
\end{definition}

\section{Hölder spaces, Sobolev spaces, and jets}

\begin{definition}[$C^k$ spaces of sections (Definition K.6)]
\label{def:chow-iv-k6-ck-sections}
\source{chow-et-al-ricci-flow-part-iv:page-0353}
\dcref{appk:K.6}
\group{Function Spaces}

For a bundle section $\phi$ with compatible connection, set
$[\phi]_j=\sup|\nabla^j\phi|$ and
$\|\phi\|_k=\sum_{j=0}^k[\phi]_j$.  The Banach space
$C^k(\mathcal M,\mathcal E)$ consists of locally $C^k$ sections with finite
$\|\cdot\|_k$ norm.
\end{definition}

\begin{definition}[Hölder spaces (Definition K.7)]
\label{def:chow-iv-k7-holder-spaces}
\source{chow-et-al-ricci-flow-part-iv:page-0354}
\uses{def:chow-iv-k6-ck-sections}
\dcref{appk:K.7}
\group{Function Spaces}


In local trivializations define
\[
[\nabla^k\phi]_\alpha=
\sup_{\beta,p\ne q\in U_\beta}
\frac{|\phi^\beta_{;i_1\cdots i_k}(p)-\phi^\beta_{;i_1\cdots i_k}(q)|}{d(p,q)^\alpha},
\qquad
\|\phi\|_{k,\alpha}=\|\phi\|_k+[\nabla^k\phi]_\alpha.
\]
The sections of finite norm form $C^{k,\alpha}(\mathcal M,\mathcal E)$.
Equivalently one may define the seminorm by parallel translation; on compact
manifolds the resulting norm is uniformly equivalent and independent of the
auxiliary choices.
\end{definition}

\begin{definition}[Sobolev spaces (Definition K.8)]
\label{def:chow-iv-k8-sobolev-spaces}
\source{chow-et-al-ricci-flow-part-iv:page-0354}
\uses{def:chow-iv-k6-ck-sections}
\dcref{appk:K.8}
\group{Function Spaces}


The Sobolev norm is
\[
\|\phi\|_{k,p}=\sum_{j=0}^k\left(\int_{\mathcal M}|\nabla^j\phi|^p\,d\mu\right)^{1/p},
\]
and $W^{k,p}$ is the completion of $C^k$ in this norm.
\end{definition}

Two pointed maps have the same $k$-jet if their coordinate derivatives through
order $k$ agree.  The fiber $J^k_{p,q}(\mathcal M,\mathcal N)$ is naturally
isomorphic to $\bigoplus_{i=1}^k(S^iT_p^*\mathcal M\otimes T_q\mathcal N)$,
or, invariantly, to the symmetrized tensors
$\operatorname{Sym}(\nabla^{i-1}dF)(p)$.  The induced Riemannian distance gives
\[
d_{C^k(\mathcal M,\mathcal N)}(F,G)=
\sup_{p\in\mathcal M}d_{J^k}\bigl([F]^k_{p,F(p)},[G]^k_{p,G(p)}\bigr),
\]
with the $C^0$ case given by the supremum target distance.  This metric space
is complete.

\begin{definition}[$C^{k,\alpha}$-closeness of hypersurfaces (Definition K.9)]
\label{def:chow-iv-k9-hypersurface-closeness}
\source{chow-et-al-ricci-flow-part-iv:page-0356,chow-et-al-ricci-flow-part-iv:page-0357}
\uses{def:chow-iv-k7-holder-spaces}
\dcref{appk:K.9}
\group{Function Spaces}


A normal graph $\mathcal Y_\varphi$ is $\varepsilon$-close to a hypersurface
$\mathcal X$ in $C^{k,\alpha}$ when
$\|\varphi\|_{C^{k,\alpha}(\mathcal X)}\le\varepsilon$.
\end{definition}

\section{Harmonic maps and their linearization}

\begin{definition}[Map energy (Definition K.10)]
\label{def:chow-iv-k10-map-energy}
\source{chow-et-al-ricci-flow-part-iv:page-0357}
\dcref{appk:K.10}
\group{Harmonic Maps}

For $f:(\mathcal M,g)\to(\mathcal N,h)$,
\[
|df|^2_{g\boxtimes h}=\operatorname{tr}_g(f^*h),\qquad
E_{g,h}(f)=\int_{\mathcal M}|df|^2_{g\boxtimes h}\,d\mu_g.
\]
\end{definition}

\begin{lemma}[Map-Laplacian in local coordinates (Lemma K.11)]
\label{lem:chow-iv-k11-map-laplacian-coordinates}
\source{chow-et-al-ricci-flow-part-iv:page-0358}
\uses{def:chow-iv-k10-map-energy}
\dcref{appk:K.11}
\group{Harmonic Maps}


The pullback connection on $T^*\mathcal M\otimes f^*T\mathcal N$ defines
$\Delta_{g,h}f=\operatorname{tr}_g(\nabla^{g\boxtimes h}df)$.  In coordinates,
\[
(\Delta_{g,h}f)^\gamma=g^{ij}\left(\partial_{ij}f^\gamma-
\Gamma(g)^k_{ij}\partial_kf^\gamma+
(\Gamma(h)^\gamma_{\alpha\beta}\circ f)\partial_if^\alpha\partial_jf^\beta\right).
\]
\end{lemma}

\begin{definition}[Harmonic map (Definition K.12)]
\label{def:chow-iv-k12-harmonic-map}
\source{chow-et-al-ricci-flow-part-iv:page-0359}
\uses{lem:chow-iv-k11-map-laplacian-coordinates}
\dcref{appk:K.12}
\group{Harmonic Maps}


A map $f:(\mathcal M,g)\to(\mathcal N,h)$ is harmonic when
$\Delta_{g,h}f=0$.
\end{definition}

\begin{example}[Minimal immersions and harmonic maps (Example K.13)]
\label{ex:chow-iv-k13-minimal-immersion}
\source{chow-et-al-ricci-flow-part-iv:page-0359}
\uses{def:chow-iv-k12-harmonic-map}
\dcref{appk:K.13}
\group{Harmonic Maps}


For a smooth immersion $f$, its mean curvature vector is
$\Delta_{f^*h,h}f$.  Hence an isometric immersion is minimal exactly when it
is harmonic.
\end{example}

\begin{lemma}[First variation of map energy (Lemma K.14)]
\label{lem:chow-iv-k14-first-variation}
\source{chow-et-al-ricci-flow-part-iv:page-0360,chow-et-al-ricci-flow-part-iv:page-0361}
\uses{def:chow-iv-k10-map-energy}
\dcref{appk:K.14}
\group{Harmonic Maps}


For a variation with field $V=\partial_sf_s|_{s=0}$,
\[
\left.\frac{d}{ds}\right|_{0}E_{g,h}(f_s)=
2\int_{\mathcal M}\langle\nabla^{f^*h}V,df\rangle\,d\mu_g
=-2\int_{\mathcal M}\langle V,\Delta_{g,h}f\rangle\,d\mu_g
\]
when $V$ is compactly supported in the interior (or under the corresponding
boundary conditions).  Thus harmonic maps are critical points.
\end{lemma}

\begin{lemma}[Second variation and index form (Lemma K.15)]
\label{lem:chow-iv-k15-second-variation}
\source{chow-et-al-ricci-flow-part-iv:page-0362,chow-et-al-ricci-flow-part-iv:page-0363,chow-et-al-ricci-flow-part-iv:page-0364}
\uses{lem:chow-iv-k14-first-variation}
\dcref{appk:K.15}
\group{Harmonic Maps}


At a harmonic map, for a compactly supported variation field $V$,
\[
\left.\frac{d^2}{ds^2}\right|_0E_{g,h}(f_s)=I(V,V)
=2\int_{\mathcal M}\left(|\nabla^{f^*h}V|^2-
\operatorname{tr}_g(R^N(V,df,df,V))\right)d\mu_g.
\]
The map is weakly stable when $I(V,V)\ge0$ for all such $V$.  Nonpositive
sectional curvature of $N$ implies weak stability.
\end{lemma}

\begin{theorem}[Existence of harmonic maps into nonpositively curved targets (Theorem K.17)]
\label{thm:chow-iv-k17-eells-sampson}
\source{chow-et-al-ricci-flow-part-iv:page-0365}
\dcref{appk:K.17}
\group{Harmonic Map Heat Flow}

If $M,N$ are closed and $N$ has nonpositive sectional curvature, every smooth
$f_0:M\to N$ is homotopic to a harmonic map.  The harmonic map heat flow
\[
\partial_tf=\Delta_{g,h}f
\]
exists for all $t\ge0$ and converges smoothly to a harmonic map; strict
negative curvature gives uniqueness in the homotopy class.
\end{theorem}

\begin{lemma}[Evolution of the energy density (Lemma K.18)]
\label{lem:chow-iv-k18-energy-density}
\source{chow-et-al-ricci-flow-part-iv:page-0365,chow-et-al-ricci-flow-part-iv:page-0366}
\uses{thm:chow-iv-k17-eells-sampson}
\dcref{appk:K.18}
\group{Harmonic Map Heat Flow}


Along the harmonic map heat flow,
\[
\partial_t|df|^2=\Delta_g|df|^2-2|\nabla df|^2
-2\operatorname{Rc}^M(df\odot df)+2R^N(df\odot df\odot df\odot df).
\]
If $\sec_N\le0$ and $\operatorname{Rc}_M\ge-Kg$, then
$\partial_t|df|^2\le\Delta_g|df|^2-2|\nabla df|^2+2K|df|^2$, hence
$|df|^2\le Ce^{2Kt}$ on a closed source.  Moreover,
\[
\partial_t|\Delta_{g,h}f|^2
=\Delta_g|\Delta_{g,h}f|^2-2|\nabla(\Delta_{g,h}f)|^2
+2R^N(\Delta_{g,h}f,df,df,\Delta_{g,h}f),
\]
so nonpositive target curvature gives a uniform maximum-principle bound.
\end{lemma}

\begin{remark}[Linearization with normal boundary condition]
\label{rem:chow-iv-k-linearization-boundary}
\source{chow-et-al-ricci-flow-part-iv:page-0367,chow-et-al-ricci-flow-part-iv:page-0368}
\dcref{appk:K.43--K.44}
\group{Harmonic Map Heat Flow}

For a boundary-preserving variation $V$ of the identity, the two components of
the linearization of the harmonic-map gauge map are
\[
\left.\partial_s\right|_0(F_s^{-1})_*(\Delta_{g,g}F_s)=\Delta_gV+\operatorname{Rc}(V),
\]
\[
\left.\partial_s\right|_0(F_s^{-1})_*((F_s)_*N)_\top
=([N,V])_\top=(\nabla_NV)_\top-\mathrm{II}(V).
\]
\end{remark}

\section{Spectrum of the Hodge--de Rham Laplacian on the sphere}

\begin{proposition}[Space of homogeneous (harmonic) polynomials (Proposition K.19)]
\label{prop:chow-iv-k19-harmonic-polynomial-decomposition}
\source{chow-et-al-ricci-flow-part-iv:page-0368,chow-et-al-ricci-flow-part-iv:page-0369}
\dcref{appk:K.19}
\group{Sphere Spectrum}

For $\mathcal P_k$ the degree-$k$ homogeneous polynomials and $\mathcal H_k$
the harmonic ones,
\[
\mathcal P_{k+2}=\mathcal H_{k+2}\oplus r^2\mathcal P_k.
\]
Consequently $\mathcal P_{2j}$ and $\mathcal P_{2j+1}$ decompose orthogonally
into the corresponding sums of $r^{2\ell}\mathcal H_{k-2\ell}$, and the
restrictions of the $\mathcal H_k$ form a complete orthogonal system on $S^n$.
\end{proposition}

\begin{proposition}[Eigenvalues of $S^n$ (Proposition K.20)]
\label{prop:chow-iv-k20-function-spectrum-sphere}
\source{chow-et-al-ricci-flow-part-iv:page-0369}
\uses{prop:chow-iv-k19-harmonic-polynomial-decomposition}
\dcref{appk:K.20}
\group{Sphere Spectrum}


The eigenvalues of the Laplacian on functions on the unit sphere are
\[
\lambda_k=k(k+n-1),\qquad k\in\mathbb N\cup\{0\},
\]
with eigenspace the restrictions of $\mathcal H_k$.
\end{proposition}

\begin{theorem}[Hodge decomposition (Theorem K.21)]
\label{thm:chow-iv-k21-hodge-decomposition}
\source{chow-et-al-ricci-flow-part-iv:page-0369,chow-et-al-ricci-flow-part-iv:page-0370}
\dcref{appk:K.21}
\group{Hodge Theory}

For a closed Riemannian manifold,
\[
\Omega^p=\Delta_d(\Omega^p)\oplus H^p
=d\delta(\Omega^p)\oplus\delta d(\Omega^p)\oplus H^p
=d(\Omega^{p-1})\oplus\delta(\Omega^{p+1})\oplus H^p,
\]
where $H^p=\ker\Delta_d$ and $\Delta_d=-(d\delta+\delta d)$ in the
convention of the source.  The operators satisfy
$\Delta_d(d\varphi)=d(\Delta_d\varphi)$ and
$\Delta_d(*\varphi)=*(\Delta_d\varphi)$.
\end{theorem}

\begin{theorem}[Closed eigenforms on $S^n$ (Theorem K.22)]
\label{thm:chow-iv-k22-closed-eigenforms}
\source{chow-et-al-ricci-flow-part-iv:page-0370}
\uses{thm:chow-iv-k21-hodge-decomposition}
\dcref{appk:K.22}
\group{Sphere Spectrum}


The $k$-th eigenvalue on closed $p$-forms on the unit sphere is
\[
\lambda_k^{[p]}(S^n)=(k+p)(n+k-p+1).
\]
The eigenspace consists of restrictions of degree-$k$ homogeneous harmonic
$p$-forms on $\mathbb R^{n+1}$.
\end{theorem}

\begin{theorem}[Co-closed eigenforms on $S^n$ (Theorem K.23)]
\label{thm:chow-iv-k23-coclosed-eigenforms}
\source{chow-et-al-ricci-flow-part-iv:page-0370}
\uses{thm:chow-iv-k22-closed-eigenforms}
\dcref{appk:K.23}
\group{Sphere Spectrum}


The $k$-th eigenvalue on co-closed $p$-forms is
\[
{}^{\mathrm{cc}}\lambda_k^{[p]}(S^n)=(k+n-p)(k+p+1).
\]
\end{theorem}

\begin{corollary}[Lowest $p$-form eigenvalue (Corollary K.24)]
\label{cor:chow-iv-k24-lowest-p-form-eigenvalue}
\source{chow-et-al-ricci-flow-part-iv:page-0370}
\uses{thm:chow-iv-k22-closed-eigenforms, thm:chow-iv-k23-coclosed-eigenforms}
\dcref{appk:K.24}
\group{Sphere Spectrum}


The lowest eigenvalue of $\Delta_d$ on $p$-forms on $S^n$ is
\[
\min\{p(n-p+1),(p+1)(n-p)\}.
\]
\end{corollary}

\begin{lemma}[Codifferential restriction to $S^n$ (Lemma K.25)]
\label{lem:chow-iv-k25-codifferential-sphere}
\source{chow-et-al-ricci-flow-part-iv:page-0371,chow-et-al-ricci-flow-part-iv:page-0372}
\dcref{appk:K.25}
\group{Sphere Spectrum}

For a $p$-form $\eta$ near $S^n\subset\mathbb R^{n+1}$ and homogeneous degree-one
extensions of tangent fields,
\[
(\delta_{S^n}\eta|_{S^n}-\delta_{\mathbb R^{n+1}}\eta)(X_1,\ldots,X_{p-1})
=(n-2p+2)\eta(\nu,X_1,\ldots,X_{p-1})
 +\nu\!\left(\eta(\nu,X_1,\ldots,X_{p-1})\right).
\]
\end{lemma}

\begin{lemma}[Hodge Laplacian restriction (Lemma K.26)]
\label{lem:chow-iv-k26-hodge-laplacian-sphere}
\source{chow-et-al-ricci-flow-part-iv:page-0372}
\uses{lem:chow-iv-k25-codifferential-sphere}
\dcref{appk:K.26}
\group{Sphere Spectrum}


If $\eta$ is a closed $p$-form on $\mathbb R^{n+1}$, then
\[
(\Delta_d^{S^n}\eta|_{S^n}-\Delta_d^{\mathbb R^{n+1}}\eta)(X_1,\ldots,X_p)
=\nu\!\left(\nu(\eta(X_1,\ldots,X_p))\right)
 +(n-2p+2)\nu(\eta(X_1,\ldots,X_p)).
\]
\end{lemma}

\begin{lemma}[Closed eigenforms from homogeneous extensions (Lemma K.27)]
\label{lem:chow-iv-k27-homogeneous-closed-forms}
\source{chow-et-al-ricci-flow-part-iv:page-0372}
\uses{lem:chow-iv-k26-hodge-laplacian-sphere}
\dcref{appk:K.27}
\group{Sphere Spectrum}


If $\eta$ is a homogeneous degree-$k$ harmonic $p$-form on $\mathbb R^{n+1}$,
then $\eta|_{S^n}$ is closed and
\[
\Delta_d^{S^n}(\eta|_{S^n})=(k+p)(n-p+k+1)\eta|_{S^n}.
\]
\end{lemma}

\begin{theorem}[Gallot--Meyer estimate (Theorem K.28)]
\label{thm:chow-iv-k28-gallot-meyer}
\source{chow-et-al-ricci-flow-part-iv:page-0373}
\uses{thm:chow-iv-k21-hodge-decomposition, cor:chow-iv-k24-lowest-p-form-eigenvalue}
\dcref{appk:K.28}
\group{Hodge Theory}


If $(\mathcal M^n,g)$ is closed and the curvature operator satisfies
$\lambda_1(\operatorname{Rm})\ge k>0$, then the first positive eigenvalue on
$p$-forms obeys
\[
\lambda_1^{[p]}(\mathcal M,g)\ge
\min\{p(n-p+1),(p+1)(n-p)\}\,k.
\]
This follows from the Bochner--Weitzenböck identity
\[
\langle\Delta_d\alpha,\alpha\rangle
=\langle\Delta\alpha,\alpha\rangle+\mathcal R_p(\alpha,\alpha)
\]
and $\mathcal R_p(\alpha,\alpha)\ge p(n-p)k|\alpha|^2$.
\end{theorem}

\begin{example}[Linearization of the map-Laplacian (Exercise K.16)]
\label{ex:chow-iv-k16-map-laplacian-linearization}
\dcref{appk:K.16}
\group{Harmonic Maps}
\source{chow-et-al-ricci-flow-part-iv:page-0364}
\uses{lem:chow-iv-k15-second-variation}
For the variation map $F(x,s)=f_s(x)$ and $V=\partial_sF|_{s=0}$, prove
\[
\nabla_{\partial/\partial s}(\Delta_{g,h}^{\mathcal M}F)
=\Delta_{g,h}V+\operatorname{tr}_g(R^N(V,df)df).
\]
\end{example}
